% Checked mathematical fragments from /root/pencil.
% Dependencies: finding environment with two arguments; amsmath, amsthm,
% amssymb; macros \BR and \Span. No division by a factorial is used.
% The final P-08 here is the finite-difference upper certificate:
% its source-note locator is PENCIL_PROOFS.md, P-09 (upper-bound part).

\section{Exact two-dimensional residual minima}

Throughout this section the field is algebraically closed and the tensors
have integral ordered-coordinate coefficients \(1\) or \(-1\). For the
polynomial families these are the divided-power coefficient tensors.
In characteristic zero they differ from the usual symmetric coefficient
tensors only by a nonzero scalar. In positive characteristic the integral
tensor, rather than evaluation of a symmetric tensor as a polynomial, is
the convention.

For a \(j\)-plane \(W\subset U_{n,b}\), choose a basis \(u_1,\ldots,u_j\)
and put \(P_W=\sum_{h=1}^j a_h\otimes u_h\). Its border rank is independent
of the basis. Write
\[
 \mu_j(U_{n,b})=\min_{W\in\operatorname{Gr}(j,U_{n,b})}\BR(P_W).
\]
The residual factors of \(P_W\) remain ungrouped.

\begin{finding}{P-01}{the Grassmannian minimum is a coordinate minimum}
In each of the four families and for every allowed \(n,b,j\),
\[
 \mu_j(U_{n,b})=
 \min_{\substack{\mathcal S\subset\mathcal B_{n,b}\\|\mathcal S|=j}}
 \BR\left(\sum_{Q\in\mathcal S}a_Q\otimes Q\right),
\]
where \(\mathcal B_{n,b}\) is the coordinate residual basis: \(Q_I\)
in a row family and \(Q_{I,J}\) in a polynomial family.
\end{finding}
\begin{proof}
For a row family, act on \(e_i\), \(0\leq i<n\), by
\(g(t)e_i=t^{2^i}e_i\). The basis tensor \(Q_I\) is an eigenvector
of weight \(\sum_{i\in I}2^i\), and these weights are distinct.
For a polynomial family act on \(e_{ij}\) by
\(t^{2^i+2^{n+j}}\). Then \(Q_{I,J}\) has weight
\(\sum_{i\in I}2^i+\sum_{j\in J}2^{n+j}\), again with distinct
weights. Integer exponents define distinct characters in every
characteristic.

Order the residual basis \(q_1,\ldots,q_N\) by increasing weights
\(w_1<\cdots<w_N\). Row reduction gives a basis of any chosen plane
of the form
\[
 u_h=q_{p_h}+\sum_{i>p_h}c_{hi}q_i,\qquad h=1,\ldots,j,
\]
with distinct pivot indices. Consequently
\[
 u_h(t):=t^{-w_{p_h}}g(t)^{\otimes b}u_h
 =q_{p_h}+\sum_{i>p_h}c_{hi}t^{w_i-w_{p_h}}q_i
\]
is polynomial in \(t\), and \(u_h(0)=q_{p_h}\).
For \(t\ne0\), the tensor \(\sum_h a_h\otimes u_h(t)\) is
obtained from \(P_W\) by invertible maps in its factors: \(g(t)\)
in each residual factor and a diagonal map in the first factor.
Its border rank therefore equals \(\BR(P_W)\).
The border-rank-at-most-\(r\) locus is closed, so specialization
at zero cannot increase border rank. This gives a coordinate plane
of no larger border rank. The converse inequality follows because
every coordinate plane is an allowed plane.
\end{proof}

\begin{finding}{P-02}{all two-dimensional minima reduce to an adjacent pair}
For \(n\geq b+1\), the minimum \(\mu_2(U_{n,b})\) is the border rank
of the following single tensor, independently of \(n\):
\[
 \begin{cases}
 a_1\otimes Q_{[b]}+
 a_2\otimes Q_{[b-1]\cup\{b+1\}},&\text{row families},\\[2mm]
 a_1\otimes Q_{[b],[b]}+
 a_2\otimes Q_{[b],[b-1]\cup\{b+1\}},&\text{polynomial families}.
 \end{cases}
\]
This is an exact smaller-tensor reduction; the displayed border
rank for arbitrary \(b>2\) is not evaluated here.
\end{finding}
\begin{proof}
By P-01 it suffices to consider coordinate pairs.
For two distinct \(b\)-subsets \(I,K\), write their common part as
\(C\), with \(|C|=c\), and their equally sized differences as
\(\{i_1,\ldots,i_d\}\) and \(\{k_1,\ldots,k_d\}\), where
\(d=b-c\geq1\). Map the common coordinates bijectively onto
\(e_1,\ldots,e_c\); map both \(e_{i_h}\) and \(e_{k_h}\) to
\(e_{c+h}\) for \(h<d\); and map \(e_{i_d},e_{k_d}\) to
\(e_b,e_{b+1}\), respectively. Send unused coordinates to zero.
The map is injective on each of the two selected coordinate sets.
Its tensor power therefore takes the two selected row tensors to
the adjacent pair, up to signs in the determinant case. Those
signs are absorbed in the first factor.

For polynomial tensors, suppose first that the column sets \(J,L\)
are distinct. Map the two row sets \(I,K\) injectively onto the
same \(b\)-element row set, by matching their differences and
retaining their intersection. Map the column sets \(J,L\) to the
adjacent pair just constructed. The induced linear map on matrix
coordinates, \(e_{ij}\mapsto e_{f(i),g(j)}\), is injective on the
rows and columns within each selected term. It therefore takes
each selected minor or permanent to the desired adjacent minor
or permanent, up to a determinant sign. If \(J=L\), then
\(I\ne K\); interchange rows and columns first, and transpose
afterwards.

Products of linear maps cannot increase border rank. Thus every
coordinate pair has border rank at least that of the adjacent
pair. The adjacent pair occurs in \(U_{n,b}\), proving equality.
\end{proof}

\begin{finding}{P-03}{a uniform star degeneration}
Let \(k\geq1\), let \(a_1,\ldots,a_k\) and \(e_0,\ldots,e_k\)
be independent, and let \(\delta\in\{1,-1\}\). Then
\[
 S_{k,\delta}:=\sum_{i=1}^k
 a_i\otimes(e_0\otimes e_i+\delta e_i\otimes e_0)
 \qquad\text{satisfies}\qquad
 \BR(S_{k,\delta})=k+1
\]
over every algebraically closed field.
\end{finding}
\begin{proof}
Contraction in the last factor gives the \(k+1\) independent
vectors
\(\delta\sum_i a_i\otimes e_i\) and \(a_i\otimes e_0\),
\(1\leq i\leq k\). Thus the corresponding flattening has rank
\(k+1\), proving the lower bound.
For \(t\ne0\), the following is a sum of \(k+1\) pure tensors:
\[
 \frac1t\left[
 \sum_{i=1}^k a_i\otimes(e_0+\delta t e_i)\otimes(e_0+t e_i)
 -\left(\sum_{i=1}^k a_i\right)\otimes e_0\otimes e_0
 \right].
\]
Its exact expansion is
\[
 S_{k,\delta}+\delta t\sum_{i=1}^k a_i\otimes e_i\otimes e_i.
\]
It therefore specializes to \(S_{k,\delta}\). There is no
division by an integer; the identity remains valid in
characteristic two.
\end{proof}

\begin{finding}{P-04}{the requested \(n=4,b=2,j=2\) computation}
For every \(n\geq3\), the exact two-dimensional residual minima are
\[
 \begin{array}{c|cccc}
 \text{family}&D_n&P_n&\mathcal D_n&\mathcal P_n\\ \hline
 \mu_2(U_{n,2})&3&3&6&6 .
 \end{array}
\]
In particular, at \(n=4\) the two-dimensional substitution
bound is \(7\) in the row families and \(40\) in the polynomial
families. These assertions hold for the integral coefficient
tensors over every algebraically closed field.
\end{finding}
\begin{proof}
The adjacent row pair in P-02 is \(S_{2,-1}\) for the
determinant and \(S_{2,1}\) for the permanent, up to labels and
signs. P-03 gives its exact border rank three.

For the polynomial pair use the six matrix-coordinate variables
\[
 a=e_{11},\ b=e_{12},\ c=e_{13},\
 d=e_{21},\ e=e_{22},\ f=e_{23}.
\]
Its two slices are the ordered coefficient tensors of
\(ae+\sigma bd\) and \(af+\sigma cd\), where \(\sigma=-1\)
for the determinant and \(+1\) for the permanent.
The associated tensor is the sum of a symmetric star
\(S_{2,1}\) on \((a,e,f)\) and \(\sigma S_{2,1}\) on
\((d,b,c)\), with the same two-dimensional first factor.
Adding the two explicit curves of P-03 gives a six-term
degeneration.

For the matching lower bound, each coordinate quadratic is a
nonsingular signed matching on its four-variable coordinate
support. Its matrix kernel is precisely the coordinate
complement of that support. For a distinct pair
\((I,J),(K,L)\), the common slice kernel is therefore the
complement of
\((I\times J)\cup(K\times L)\). The pencil flattening has rank
\[
 |(I\times J)\cup(K\times L)|
   =8-|I\cap K|\,|J\cap L|\geq6.
\]
The last inequality uses distinctness of the pairs; the only
excluded product larger than two is \(2\cdot2\), which would
make the pairs identical. The adjacent pair attains six.
P-01 proves uniformity over the entire Grassmannian, not only
coordinate planes.

Finally the generalized substitution theorem R-01 gives
\(\BR(T_n)\geq\binom n2+1\) for the row families and
\(\BR(T_n)\geq\binom n2^2+4\) for the polynomial families.
At \(n=4\) these are the stated numbers.
\end{proof}

\begin{finding}{P-05}{exact Koszul certificates for the pencils}
For a pencil \(P=a_1\otimes Q_1+a_2\otimes Q_2\) in
\(A\otimes B\otimes C\), with \(\dim A=2\), the \(p=1\)
Koszul map
\[
 K_P:A\otimes B^*\longrightarrow\bigwedge^2A\otimes C,
 \qquad
 a\otimes\beta\longmapsto
 \sum_{i=1}^2a_i\wedge a\otimes Q_i(\beta)
\]
has matrix \([-Q_2\mid Q_1]\), up to an overall sign convention.
Its pure-tensor cost is \(\binom11=1\).
For the coordinate row pencils its rank is the union size of
the two edges; for the coordinate polynomial pencils its rank
is \(8-|I\cap K||J\cap L|\). Thus P-04 is certified by
Koszul flattenings with exact costs.
\end{finding}
\begin{proof}
For a pure tensor the displayed map has rank at most one.
Rank subadditivity and closed determinantal rank loci give
\(\BR(P)\geq\operatorname{rank}K_P\).
The image of each slice is its coordinate support, proving
the stated concatenated ranks. More explicitly, for each
variable in the support union choose a slice containing it
and the unique input coordinate paired with it in that
signed matching. These are distinct columns, each with one
nonzero entry in the chosen rows, and give a square minor
with determinant \(1\) or \(-1\).

At \(n=4\), exhaustive integer certificates give \(15\) row
pairs, with rank distribution \(12\) of rank three and \(3\)
of rank four, and \(630\) polynomial pairs, with \(144\) of
rank six, \(288\) of rank seven, and \(198\) of rank eight.
The checker \path{check_pencils.py} constructs the minors
and verifies both signs, as well as the star degenerations.
The two-dimensional Koszul map is an ordinary flattening
after a basis change; no stronger Koszul estimate is needed
for this exact minimum.
\end{proof}

\section{Larger residual subspaces}

\begin{finding}{P-06}{anchored collections have explicit exact border rank}
Let \(H\subset\{1,\ldots,r\}\times\{1,\ldots,s\}\) use every
row and column vertex, let \(h=|H|>0\), and let
\(a,x_i,y_j,z_{ij}\), \((i,j)\in H\), be independent.
For \(\sigma=\pm1\), define
\[
 T_H=\sum_{(i,j)\in H}c_{ij}\otimes
 \bigl(a\otimes z_{ij}+z_{ij}\otimes a
       +\sigma(x_i\otimes y_j+y_j\otimes x_i)\bigr).
\]
Then \(\BR(T_H)=h+r+s+1\), in every characteristic.
\end{finding}
\begin{proof}
Every slice is a nonsingular signed matching on its
four-variable support. The union contains all
\(h+r+s+1\) variables, so the last-factor flattening has
that rank.

For the upper bound put
\[
 \begin{split}
 v_{00}&=a,\qquad v_{i0}=a+t x_i,\qquad v_{0j}=a+t y_j,\\
 v_{ij}&=a+t x_i+t y_j+\sigma t^2z_{ij}.
 \end{split}
\]
The curve
\[
 T(t)=\sigma t^{-2}\sum_{(i,j)\in H}c_{ij}\otimes
 \bigl(v_{ij}^{\otimes2}-v_{i0}^{\otimes2}
       -v_{0j}^{\otimes2}+v_{00}^{\otimes2}\bigr)
\]
is polynomial after cancellation and has value \(T_H\) at
zero. Indeed, its constant term after division by \(t^2\)
is exactly the displayed slice. Grouping by the \(h+r+s+1\)
vectors \(v_{ij},v_{i0},v_{0j},v_{00}\) expresses it as at
most that many pure tensors for \(t\ne0\).
The identity is over \(\mathbb Z\).
\end{proof}

\begin{finding}{P-07}{determinant coordinate minima reduce to order ideals}
For either determinant family, P-01 can be restricted to
coordinate subsets that are lower order ideals of cardinality
\(j\). The order is componentwise order on increasing
\(b\)-subsets in the row family, and product componentwise
order on pairs of increasing \(b\)-subsets in the polynomial
family. The assertion holds in every characteristic.
\end{finding}
\begin{proof}
Fix \(v>u\) and apply the transvection
\(e_v\mapsto e_v+t e_u\). In the exterior-power coordinate
basis, it sends
\[
 Q_I\longmapsto Q_I+\varepsilon_I t
 Q_{(I\setminus\{v\})\cup\{u\}}
\]
when \(v\in I\), \(u\notin I\), and fixes the vector
otherwise. If both indices occur, the added exterior
product has a repeated coordinate and vanishes. Thus the
nontrivial actions form disjoint two-element chains
with fixed lower partner. For polynomial minors apply
the same construction separately to row or column
exterior powers.

If both partners of a chain are selected, their span is
unchanged. If only the lower partner is selected, it is
fixed. If only the upper partner is selected, substitute
\(t=s^{-1}\) and normalize its frame vector to
\(sQ_{\rm upper}+\varepsilon Q_{\rm lower}\); it tends to
the lower partner. Different chains have different
partners, so the limiting frame remains independent.
For \(s\ne0\) the residual tensors differ only by product
general linear transformations and changes of first-factor
basis. Border rank cannot increase at the limit.

Whenever a replacement occurs, the sum of all selected
integer coordinate indices strictly decreases. Iterating
must terminate. At termination the selected subset is
closed under replacement of any index by a smaller absent
index. Such replacements generate componentwise order:
to pass from \(I\) to \(J\leq I\), successively replace the
first unequal coordinate by that of \(J\). The polynomial
case gives the product order by doing rows and columns
separately. The resulting coordinate plane is therefore
an order ideal and has no larger border rank.
\end{proof}

The exterior-power cancellation in P-07 is essential.
The analogous permanent image is not invariant under
these transvections, so P-07 is not asserted for it.

\begin{finding}{P-08}{finite differences for the exceptional permanent subspaces}
Let \(T_5\) have the five ordered coefficient slices
\[
 ae+bd,\quad af+cd,\quad bf+ce,\quad ah+bg,\quad ai+cg
\]
in the row-major variables
\(\left(\begin{smallmatrix}a&b&c\\d&e&f\\g&h&i\end{smallmatrix}\right)\).
Then \(\BR(T_5)\leq11\).
Adding the slice \(bi+ch\) gives \(T_6\) with
\(\BR(T_6)\leq12\).
Both upper bounds have integral, explicit finite-difference
certificates. In conjunction with R-04, \(\BR(T_5)=11\).
\end{finding}
\begin{proof}
Put
\[
 v(u,w,y,z)=(r_0+u r_1+w r_2)\otimes(s_0+y s_1+z s_2),
 \qquad F=v\otimes v.
\]
The coefficients of \(uy,uz,uyz,wy,wz\) in \(F\) are
exactly the five stated ordered symmetric quadrics.
For example, the coefficient of \(uy\) is
\(a\otimes e+e\otimes a+b\otimes d+d\otimes b\);
that of \(uyz\) is
\(b\otimes f+f\otimes b+c\otimes e+e\otimes c\).

For a set \(S\) of parameters define
\[
 D_S(t)=t^{-|S|}
 \sum_{R\subset S}(-1)^{|S|-|R|}F(t\,1_R).
\]
Inclusion-exclusion kills every monomial with support
properly contained in \(S\); parameters outside \(S\)
are zero. After cancellation \(D_S(t)\) is polynomial
and its value at zero is the coefficient of the
squarefree monomial supported on \(S\).

Take the first-factor-labelled sum of these differences
for \(S=uy,uz,uyz,wy,wz\).
The distinct evaluation sets are the eight subsets of
\(\{u,y,z\}\) and the three additional sets
\(\{w\},\{w,y\},\{w,z\}\). Grouping their first-factor
coefficients yields eleven pure tensors, with the desired
limit. For \(T_6\), the additional set \(S=wyz\) adds only
the evaluation set \(\{w,y,z\}\), giving twelve terms.
The checker \path{check_finite_differences.py} expands
the Laurent expressions over \(\mathbb Z\), verifies
cancellation of all negative powers, and checks every
constant tensor coefficient.
\end{proof}

\subsection*{Scheme and representation dictionary for P-01--P-08}
The row and column diagonal actions are split tori inside
general linear group schemes over \(\mathbb Z\). In P-01 the
echelon frame is a rank-\(j\) subbundle of the trivial residual
bundle on \(\mathbb A^1\), because its pivot coefficient
submatrix remains the identity. Its tautological tensor
specializes in the closed Segre secant cone.
P-02 uses morphisms of free modules, hence of their tensor
bundles. P-03, P-06 and P-08 supply integral curves in
secant cones. P-05 is an exterior-power morphism of locally
free sheaves; its unit minors give the asserted
determinantal inequalities after every base change.
In P-07 the elementary lowering root subgroups and the
diagonal torus generate a split Borel; the terminal
coordinate planes are its invariant subspaces.
This is an explicit form of the projective Borel-fixed
reduction, requiring no semisimplicity in positive
characteristic. The exact EGA, SGA and FGA locators for
these constructions are supplied in the dictionary
section of the main paper.
